%\subsection{Training Pipeline}
%\label{sec:training}

%This section describes the training pipeline for DeXposure-FM, including data construction, multi-task objectives, and optimization. All graph and flow notation follows Section~\ref{sec:formalisation} (credit-exposure definition in Eq.~\eqref{eq:credit-exposure-definition}, node weights in Eq.~\eqref{eq:node-weight}, token-level flows in Eq.~\eqref{eq:value-flow}, and edge weights in Eq.~\eqref{eq:edge-weight}).

\subsection{Training Data Preparation}

For each discrete interval $\tau=t_2-t_1$ we construct a weighted directed snapshot $G_{\tau}=(\mathcal P_{\tau},\mathcal E_{\tau})$ with node weights $w_{\tau}(p)$ and edge weights $w_{\tau}(e_{pq})$. We form rolling training examples indexed by an anchor time $\tau$:
\begin{equation}
\underbrace{\{G_{\tau-L+1},\ldots,G_{\tau}\}}_{\text{context window}}
\;\longrightarrow\;
\underbrace{\{G_{\tau+1},\ldots,G_{\tau+H}\}}_{\text{forecast targets}},
\end{equation}
where $L$ is the context length and $H$ is the forecast horizon. For each protocol $p\in\mathcal P_{\tau}$ we provide node features (log-transformed TVL, token count, composition entropy, and protocol category) along with the graph structure $G_{\tau}$.

\paragraph{Time-based rolling splits.}
We evaluate models using time-based rolling splits (walk-forward validation), ensuring that all validation and test targets occur strictly after the corresponding training window to avoid look-ahead bias.

\paragraph{Scalable edge sampling.}
Since $|\mathcal E_{\tau}|$ can be large, we train with edge sampling. For each anchor $\tau$ and horizon $h$, we draw a set of positive edges from $\mathcal E_{\tau+h}$ with a mixture distribution that upweights large exposures (proportional to $w_{\tau+h}(e_{pq})$) while retaining a uniform component to preserve coverage of smaller links. When training an auxiliary link-existence head, we additionally sample non-edges as negatives.

\subsection{Training Pipeline}
This section describes the training pipeline.

\subsubsection{Training phases}
We use a two-phase training schedule:

\paragraph{Phase A: frozen encoder (linear probing).}
We initialize the GraphPFN encoder from a pre-trained checkpoint and freeze all encoder weights. We train only the task-specific heads (link predictor and node predictor) until validation loss stabilizes. This phase evaluates the quality of pre-trained representations.

\paragraph{Phase B: fine-tuned encoder.}
We unfreeze all encoder parameters and train end-to-end with a unified learning rate. This allows the model to adapt representations to the DeXposure domain while leveraging pre-trained initialization.

\subsubsection{Loss functions}
\label{sec:training:loss}
We train DeXposure-FM with a weighted multi-task objective that reflects heavy-tailed magnitudes and the economic importance of large exposures.

\paragraph{Edge forecasting loss.}
We predict future edge weights $\widehat w_{\tau+h}(e_{pq})$ for horizons $h \in \{1, 3, 7\}$ weeks. Because exposures are heavy-tailed, we train in log space:
\begin{equation}
y_{\tau+h}(e_{pq}) = \log\bigl(1+w_{\tau+h}(e_{pq})\bigr),\qquad
\hat y_{\tau+h}(e_{pq}) = \log\bigl(1+\widehat w_{\tau+h}(e_{pq})\bigr).
\end{equation}
We use a robust regression loss (SmoothL1/Huber):
\begin{equation}
\mathcal L_{\mathrm{edge}} = \mathbb E_{(p,q)\in \mathcal S^+_{\tau+h}} \left[ \mathrm{SmoothL1}\!\left(\hat y_{\tau+h}(e_{pq})-y_{\tau+h}(e_{pq})\right) \right],
\label{eq:loss-edge}
\end{equation}
where $\mathcal S^+_{\tau+h}$ is the set of positive (existing) edges at horizon $h$.

\paragraph{Log-space transformation.}
To handle heavy-tailed exposure magnitudes, we train in log space. The log-transform naturally downweights extreme values and improves gradient stability.

\paragraph{Link existence loss.}
To handle network sparsity, we decompose forecasting into edge existence prediction and conditional weight prediction. The existence head is trained with binary cross entropy:
\begin{equation}
\mathcal L_{\mathrm{link}} = \mathbb E_{(p,q)\in \mathcal S_{\tau+h}} \Bigl[\mathrm{BCE}\bigl(\hat a_{\tau+h}(p,q), a_{\tau+h}(p,q)\bigr)\Bigr],
\label{eq:loss-link}
\end{equation}
where $\mathcal S_{\tau+h}$ includes both positive edges and negative samples (non-edges).

\paragraph{Node forecasting loss.}
We forecast node TVL using the same log transform:
\begin{equation}
\mathcal L_{\mathrm{node}} = \mathbb E_{p\in \mathcal P_{\tau+h}} \left[ \mathrm{SmoothL1}\!\left(\log(1+\widehat w_{\tau+h}(p))-\log(1+w_{\tau+h}(p))\right) \right].
\label{eq:loss-node}
\end{equation}

\paragraph{Imputation evaluation.}
We evaluate the model's imputation capability by randomly masking edges at test time and measuring reconstruction accuracy. This tests the model's ability to infer missing exposure information from network context.

\paragraph{Shock analysis evaluation.}
We evaluate model robustness during historical stress events (Terra/Luna collapse in May 2022, FTX collapse in November 2022) by measuring performance degradation compared to normal market conditions. Network-level statistics (Gini coefficient, HHI, density) are computed from predicted graphs for analysis.

\paragraph{Overall objective.}
The training objective combines edge existence, edge weight, and node prediction losses:
\begin{equation}
\mathcal L = \lambda_{\mathrm{link}}\mathcal L_{\mathrm{link}} + \lambda_{\mathrm{edge}}\mathcal L_{\mathrm{edge}} + \lambda_{\mathrm{node}}\mathcal L_{\mathrm{node}}.
\label{eq:loss-total}
\end{equation}
We tune $\{\lambda_{\cdot}\}$ on validation data; default values are $\lambda_{\mathrm{link}}=1$, $\lambda_{\mathrm{edge}}=1$, $\lambda_{\mathrm{node}}=0.5$.

\subsubsection{Optimizer and learning-rate schedule}
We optimize with Adam and use early stopping based on validation loss.

\paragraph{Base optimizer.}
We use Adam with learning rate $\mathrm{lr} \approx 10^{-3}$ and weight decay $10^{-4}$.

\paragraph{Training configuration.}
Default hyperparameters include: negative sampling ratio 5:1, 5--10 training epochs, and batch processing of edge pairs for memory efficiency. We use early stopping based on validation AUPRC.

\begin{table}[t]
\centering
\small
\setlength{\tabcolsep}{6pt}
\renewcommand{\arraystretch}{1.05}
\begin{tabular}{p{0.46\linewidth} p{0.46\linewidth}}
\hline
\textbf{Hyperparameter} & \textbf{Value / tuning range} \\
\hline
Time step & Weekly intervals $\tau$ \\
Forecast horizon $h$ & $\{1, 3, 7\}$ weeks \\
Targets & $\log(1+w_{\tau}(p))$ and $\log(1+w_{\tau}(e_{pq}))$ \\
Robust loss & SmoothL1 (Huber) \\
Edge sampling & neg:pos $= 5:1$ \; (tuning: 3--10) \\
Optimizer & Adam; lr $= 10^{-3}$; wd $= 10^{-4}$ \\
Hidden dimension & $256$ \; (tuning: $\{128, 256, 512\}$) \\
Training epochs & $5$--$15$ \\
Early stopping & Based on validation AUPRC \\
Loss weights & $\lambda_{\mathrm{link}}=1$, $\lambda_{\mathrm{edge}}=1$, $\lambda_{\mathrm{node}}=0.5$ \\
\hline
\end{tabular}
\caption{Training hyperparameters. Hyperparameter tuning is performed using Optuna with TPE sampling and median pruning.}
\label{tab:training-hparams}
\end{table}
